\section{Related Work}
Literature on MIAs initially focused on discriminative models \citep{shokri2017membership} but only recently started to be applied to synthetic data generators \citep{stadler2022synthetic, hu2022membership}.
Early methods approximated the generator's output density to detect overfitting. \citet{hilprecht2019monte} proposed a Monte Carlo attack counting generated points near a target record, while \citet{chen2020gan} introduced GAN-Leaks, using the distance to the k-nearest synthetic neighbors as a density surrogate. \citet{hayes2019logan} introduced LOGAN, leveraging the generator's discriminator or a simple classifier to score records.

smMIAs were first introduced in \citet{shokri2017membership} as a powerful attack against machine learning models, and later adapted to synthetic data generators in \citet{stadler2022synthetic}, where the authors show that synthetic data generators are vulnerable to MIAs and that the success rate is correlated with the quality of the synthetic data. In this work, feature extraction is based on summary statistics, correlations, and marginal histograms of shadow synthetic datasets.
This framework was improved in \citet{houssiau2022tapas}, where the extracted features are the results of counting queries on the shadow synthetic datasets. In particular, the authors use k-way marginal statistics computed over subsets of the attribute values of the targeted record.
Building on this, \citet{meeus2023achilles} show how the success rate of smMIAs can be significantly improved by selecting outliers as target records.
Furthermore, \citet{guepin2023synthetic} relaxed the common requirement of a reference dataset, at the cost of a significant decrease in the attack's efficiency and efficacy.

In \citet{van2023membership}, the authors develop DOMIAS, a density-based MIA that exploits the local overfitting of the generative model to infer membership. DOMIAS also requires access to an auxiliary dataset, but not to the generative model, and it has been shown to be more effective than many MIA methods, while being more computationally efficient than smMIAs. Despite being an efficient way to measure privacy risk, DOMIAS has not been directly compared to smMIAs.
